% !TeX root = ../main.tex
% ---------------------------------------------------------------------------
% Engineering Geology for Underground Works
% Scope extension: not derived from the RSE3010 weekly material.
% Key sources: Hoek, Practical Rock Engineering; Waltham, Foundations of
% Engineering Geology; Knill (2003) Rankine Chapter on the ground model.
% ---------------------------------------------------------------------------
\chapter{Engineering Geology for Underground Works}
\label{ch:engineering-geology}


\section{The geological model and the ground model}
\label{sec:ground-model}
\scopenote{The distinction between a geological interpretation, an engineering
ground model with parameters attached, and the idealised design model used in
calculation. Knill's framing of the ground model as the central engineering
deliverable.}

\section{Rock materials and weathering}
\label{sec:rock-materials-and-weathering}

\subsection{Rock-forming processes and engineering behaviour}
\label{sec:rock-types}
\scopenote{Igneous, sedimentary and metamorphic rocks organised by what they
do in an excavation: anisotropy in schists and shales, bedding control in
sedimentary sequences, jointing patterns in intrusives.}

\subsection{Weathering profiles and residual soils}
\label{sec:weathering}
\scopenote{Weathering grade classification, the rock-soil transition, and why
the weathering profile often governs portal and shallow-tunnel design. The
Anderson Creek Formation profile in Chapter~\ref{ch:site-investigation} as the
worked local example.}

\section{Discontinuities and structure}
\label{sec:discontinuities-and-structure}

\subsection{Discontinuities: origin, geometry and description}
\label{sec:discontinuities}
\scopenote{Joint sets, bedding, foliation, faults. The ISRM description
parameters --- orientation, spacing, persistence, roughness, aperture,
infilling, seepage --- and which of them each design method consumes.}

\subsection{Structural geology along an alignment}
\label{sec:structural-geology}
\scopenote{Folds, faults and shear zones as alignment-scale features. Fault
zone architecture (core, damage zone), and why fault crossings dominate risk
registers out of proportion to their length.}

\subsection{Discontinuity strength and behaviour}
\label{sec:discontinuity-strength}
\scopenote{Patton's bilinear model, Barton--Bandis JRC--JCS, scale effects on
roughness, and infill-controlled strength. This underpins the wedge analyses
in Chapters~\ref{ch:caverns} and~\ref{ch:rock-reinforcement}.}

\section{Stress history and problem ground}
\label{sec:stress-history-and-problem-ground}

\subsection{Stress history and locked-in stress}
\label{sec:stress-history}
\scopenote{Overconsolidation, erosional unloading and residual stress as
sources of the horizontal stresses quantified in
Chapter~\ref{ch:in-situ-stress}.}

\subsection{Problem ground}
\label{sec:problem-ground}
\scopenote{Karst and solution features; swelling clays and anhydrite; gas-bearing
and coal-measure strata; ultra-weak and heterogeneous rock masses such as
flysch. Geological indicators that should trigger further investigation.}

\section{Hazards, uncertainty and baselining}
\label{sec:hazards-uncertainty-and-baselining}

\subsection{Geological hazards and the alignment}
\label{sec:geological-hazards}
\scopenote{Hazard mapping along a route, and how alignment and depth selection
are themselves geotechnical design decisions.}

\subsection{Geological uncertainty and baselining}
\label{sec:geological-baselining}
\scopenote{Characterising what is not known, and how the Geotechnical Baseline
Report converts geological uncertainty into contractual risk allocation. Full
treatment in Chapter~\ref{ch:risk-delivery}.}
